\chapter{Choice Theory}

The utility-maximization approach to choice has several features that
help explain its long and continuing dominance in economic analysis:

\begin{itemize}
\item
  \textbf{Normative usefulness.} It ties individual choices to a welfare
  criterion that can be compared with the government's, which
  underwrites both policy analysis and the modern democratic emphasis on
  respecting individual preferences.
\item
  \textbf{Positive predictions.} It yields sharp comparative-statics
  predictions: how choices respond when prices, incomes, or other
  parameters change.
\item
  \textbf{Wide scope.} The same machinery applies across consumption,
  labor supply, finance, and beyond.
\item
  \textbf{Compactness.} Empirical predictions follow from a sparse
  model: a description of the chooser's objectives and constraints,
  nothing more.
\end{itemize}

\section{Preference-Based Approach}\label{preference-based-approach}

Rational choice theory starts from the idea that individuals have
preferences and choose so as to maximize utility, where the utility is
itself a representation of those preferences. Our primary task is to
formalize what that statement means and to extract from it precise
predictions about the pattern of decision-making we should observe.

\subsection{Preference Relation}\label{preference-relation}

\defn{(Weak) Preference Relation}{
Let \(X\) be a set of possible choices. Consider a \textit{weak preference relation} $\pref$ over the
set \(X\), as a binary relationship:
\[x \succeq y \iff \text{``}x \text{ is at least as good as }y \text{"}\]
}

\rmkb{
The weak preference relation \(\succeq\) implies the associated strict
preference relation \(\succ\) and indifference relation \(\sim\):

\begin{itemize}
\item
  \(x\) is \textit{strictly preferred} to \(y\), or \(x \succ y\), if
  \(x\succeq y\) but not \(y \succeq x\).
\item
  \(x\) is \textit{indifferent} to \(y\), or \(x \sim y\), if \(x\succeq y\) and
  \(y \succeq x\).
\end{itemize}
}

With the weak preference relation in hand, we now impose two assumptions
that together formalize what we mean by a \emph{rational} preference.

The first is \emph{completeness}: an agent is never undecided when
faced with two choices.

\defn{Completeness}{
A preference relation \(\succeq\) on
\(X\) is \emph{complete}, if for all \(x, y \in X\), either
\(x \succeq y\), or \(y \succeq x\), or both.
}

The second is \emph{transitivity}: an agent's weak preference cannot
cycle, except among choices to which the agent is indifferent.

\defn{Transitivity}{
A preference relation \(\succeq\) on
\(X\) is \emph{{transitive}}, if whenever \(x \succeq y\) and
\(y\succeq z\), \(x \succeq z\).
}

From completeness and transitivity, we have the following corollaries:

\cor{}{
\begin{enumerate}
\item
  While the definition of transitivity involves only 3-choice cycles, it also
  extends to all \(n\)-object cycles, i.e., that it implies that for any \(n\)
  choices \(x_1,x_2,\cdots,x_n\in X\) such that \(x_1 \succeq x_2\),
  \(x_2\succeq x_3\), ..., \(x_{n-1} \succeq x_n\), we must have
  \(x_1 \succeq x_n\). (Hint: use induction on \(n\).)
\item
  Transitivity of a weak preference relation \(\succeq\) implies
  transitivity of the associated strict preference relation \(\succ\) and
  the indifference relation \(\sim\).
\end{enumerate}
}

\rmkb{
\begin{enumerate}
\item
  We take the \emph{weak} preference relation \(\pref\) as primitive
  (rather than the strict relation \(\succ\)) because completeness is
  the natural axiom for \(\pref\): with strict preference, we would need
  a third case for indifference, which is messier.
\item
  Transitivity is inconsistent with certain ``framing effects'' that
  show up in experimental data — e.g., when the order in which options
  are presented systematically changes the chosen alternative.
\end{enumerate}
}

Completeness and transitivity together let us formalize rationality.

\defn{Rationality}{
A preference relation \(\pref\) on
\(X\) is \emph{rational} if it is both \emph{complete} and
\emph{transitive}.
}

\subsection{Choice Rule}\label{choice-rule}

Given preferences, the agent's \emph{choice rule} on an ``opportunity
set'' \(B\se X\) induced by the preference relation \(\succeq\) is
\[C_{\succeq}(B) = \{x\in B \mid x\succeq y,\;\forall y\in B\},\]
i.e., the set of items in \(B\) the agent likes at least as much as
every other alternative in \(B\) — the most-preferred options.

\rmkb{
\begin{enumerate}

\item
  \(C_\succeq(B)\) may contain more than one element — ties between
  equally-preferred alternatives are perfectly admissible.
\item
  \(C_\succeq(B)\) can also be \emph{empty}. The simplest sufficient
  condition for non-emptiness is finiteness of \(B\): if \(B\) is finite
  and non-empty, then \(C_{\succeq}(B)\) is non-empty.

  \begin{itemize}
  \item
    For an example where \(C_{\pref}(B)\) is empty, take the preference
    relation \(x\succeq y \iff x\ge y\), with \(X = [0,+\infty)\) and
    \(B = (0,1)\). The supremum \(1\) is not in \(B\), so no element of
    \(B\) is most preferred.
  \item
    For infinite choice sets, we will later add technical assumptions
    (compactness of the choice set, continuity of the preference
    relation) that guarantee a choice exists.
  \end{itemize}
\end{enumerate}
}

\propp{}{
Suppose \(\succeq\) is complete and
transitive. Then, for every finite non-empty set \(B\),
\[C_\succeq(B) \ne \varnothing.\]
}{
Proceed by mathematical induction
on the number of elements of \(B\).

\begin{itemize}
\item
  \(|B| = 1\), say \(B = \{x\}\).

  \begin{itemize}
  \item
    By completeness, \(x\succeq x\), so \(x \in C_\succeq(B)\).
    \(C_\succeq(B)\ne\varnothing,\forall |B| = 1\).
  \end{itemize}
\item
  Fix \(n\ge 1\) and suppose that for all sets \(B\) with exactly \(n\)
  elements, \(C_\succeq (B)\ne \varnothing\). Next we move on to examine the case of \(|B| = n+1\).

  \begin{itemize}
  \item
    Take any \(B_n\) with \(|B_n| = n\). By the inductive hypothesis
    \(C_\succeq(B_n)\ne \varnothing\), so pick some
    \(x^*\in C_\succeq(B_n)\). Let \(B = B_n\cup \{x_{n+1}\}\).
  \item
    By completeness, we have only two (not mutually-exclusive)
    possibilities:

    \begin{itemize}
    \item
      If \(x^*\succeq x_{n+1}\), then by definition
      \(x^*\in C_\succeq(B)\), so \(C_{\succeq}(B)\ne\varnothing\).
    \item
      If \(x_{n+1}\succeq x^*\). Since \(x^*\in C_{\succeq}(B_n)\), by
      definition, \(x^*\succeq y,\forall y \in B_n\). By transitivity,
      this implies \(x_{n+1}\succeq y,\forall y\in B\). Therefore,
      \(x_{n+1}\in C_\succeq (B)\), so \(C_\succeq (B)\ne \varnothing\).
    \end{itemize}
  \end{itemize}
\item
  Hence, for every set \(B\) with exact \(n+1\) elements,
  \(C_\succeq(B)\ne\varnothing\). By the principle of mathematical induction,
  it follows that for every finite set \(B\) that is non-empty, \(C_\succeq(B)\ne \varnothing\).
\end{itemize}
}

\section{Choice-Based Approach}\label{choice-based-approach}

Much empirical work runs in the opposite direction. Rather than starting
from preferences and predicting choices, it observes choices and tries
to \emph{rationalize} them: to determine whether the observed choices
are compatible with preference maximization and, if so, what they imply
about the underlying preferences. Under this approach, the choice rule
is the primitive object of the theory, and preferences (if they exist)
are derived from it.

\defn{Choice Rule}{
Let \(\cal B\) be the set of all
nonempty subsets of \(X\)
(\(\B = 2^X\backslash \varnothing = \{B\ne \varnothing: B \subset X\}\)). A
\emph{{choice rule}} is a function \(C: \cal B\to \cal B\) with
the property that for all \(B\in \cal B\), \(C(B)\se B\).
}

\rmkb{
\begin{itemize}
\item
  \(\cal B\) is the set of all nonempty \textbf{subsets} of \(X\), which
  means all possible set of available choice(s) the agent is facing. And
  the choice rule \(C\) is a mapping from \(\cal B\) to \(\cal B\),
  which means that the agent is choosing from a set of available
  choice(s) to pick his most preferred choice(s), also a subset of
  \(X\).
\item
  Here we assume that we can see the agent choose from \emph{all}
  possible subsets of \(X\), and that the agent reports \emph{all} of
  his optimal choices from a given opportunity set.
\end{itemize}
}

The bare definition imposes no structure on either the choice rule or
any preference relation behind it. Two questions arise:

\begin{itemize}
\item
  If the rule does come from maximizing some underlying preferences,
  what can we infer about those preferences from the rule alone?
\item
  Is the rule consistent with the maximization of \emph{some} complete
  and transitive preference relation — i.e., is it
  \emph{rationalizable}?
\end{itemize}

Consider the first question. Suppose the choice rule \(C\) is consistent
with maximization of some preference relation \(\pref\), i.e.,
\(C(\cdot) = C_\pref(\cdot)\). Then, observing for some \(A\se X\) that
\(y\in A\) and \(x\in C(A)\) — \(x\) is chosen when \(y\) is available
— lets us infer \(x\pref y\). This in turn implies that for any
\(B\se X\) with \(x\in B\) and \(y\in C(B)\), we must also have
\(x\in C(B)\): we know \(x\pref y\) and \(y\pref z\) for all
\(z\in B\), so by transitivity \(x\pref z\) for all \(z\in B\). By a
symmetric argument we also get \(y\in C(A)\). So any rationalizable
choice rule must satisfy the following property as a \textit{necessary}
condition:

\defn{HARP}{
A choice function \(C: \B \to \B\) satisfies
\emph{{Houthaker's Axiom of Revealed Preference
(HARP)}} if, whenever \(x,y\in A\cap B\), and \(x\in C(A)\) and
\(y\in C(B)\), we have \(x\in C(B)\) and \(y\in C(A)\).
}


In words: if \(x\) and \(y\) are both available in two different choice
problems, and \(x\) is chosen from one while \(y\) is chosen from the
other, then \(x\) and \(y\) must \emph{both} be chosen in \emph{both}
problems. HARP rules out the obvious form of inconsistency in observed
choices — picking \(x\) over \(y\) here and \(y\) over \(x\) there.

We have just shown HARP is necessary for rationalizability. The
striking result is that it is also \emph{sufficient}:

\propp{}{
Suppose \(C: \B \to \B\) is nonempty-valued.
Then there exists a rational (complete and transitive) preference
relation \(\pref\) on \(X\) such that \(C(\cdot) = C_\pref(\cdot)\) if
and only if \(C\) satisfies HARP.
}{
\begin{itemize}
\item
  ``Only if'' part: shown above.
\item
  ``If'' part. Suppose \(C\) satisfies HARP. We construct a rational
  preference relation \(\pref_C\) that rationalizes \(C\).

  \begin{itemize}
  \item
    \textit{Define the revealed preference relation.} Set
    \(x\pref_C y\) iff there exists some \(A\se X\) with \(y\in A\) and
    \(x\in C(A)\). That is, \(x\) is revealed preferred to \(y\)
    whenever the agent picks \(x\) from some menu that also contained
    \(y\).
  \item
    \textit{\(\pref_C\) is complete.} For any \(x,y\in X\), the set
    \(C(\{x,y\})\) is non-empty, so it contains \(x\), \(y\), or both.
    In each case the corresponding ranking \(x\pref_C y\) or
    \(y\pref_C x\) follows by construction.
  \item
    \textit{\(\pref_C\) is transitive.} Suppose \(x\pref_C y\) and
    \(y\pref_C z\). The set \(C(\{x,y,z\})\) is non-empty, so it
    contains at least one of \(x,y,z\) (the cases are not mutually
    exclusive, but we just need one):

      \begin{itemize}
      \item
        \(x\in C(\{x,y,z\})\). Since \(z\in\{x,y,z\}\), the definition
        of \(\pref_C\) gives \(x\pref_C z\) directly.
      \item
        \(y\in C(\{x,y,z\})\). From \(x\pref_C y\) there exists \(A_0\)
        with \(y\in A_0\) and \(x\in C(A_0)\). Apply HARP with this
        \(A_0\) and \(B=\{x,y,z\}\) (both contain \(x\) and \(y\)) to
        get \(x\in C(\{x,y,z\})\). This reduces to case 1, so
        \(x\pref_C z\).
      \item
        \(z\in C(\{x,y,z\})\). From \(y\pref_C z\) there exists \(A_0\)
        with \(z\in A_0\) and \(y\in C(A_0)\). HARP gives
        \(y\in C(\{x,y,z\})\), reducing to case 2, and hence
        \(x\pref_C z\).
      \end{itemize}
  \item
    \textit{\(C(\cdot) = C_{\pref_C}(\cdot)\).} It suffices to show that
    for all \(A\in\B\) and \(x\in A\),
    \(x\in C(A) \iff x\pref_C y\) for all \(y\in A\).

      \begin{itemize}
      \item
        ``\(\Rightarrow\)'' holds directly by the definition of
        \(\pref_C\).
      \item
        ``\(\Leftarrow\)'' Take any \(x\in C_{\pref_C}(A)\), so
        \(x\pref_C y\) for all \(y\in A\). Since \(C(A)\) is non-empty,
        there is some \(y_0\in C(A)\). The definition of \(\pref_C\)
        gives an \(A_0\) with \(x,y_0\in A_0\) and \(x\in C(A_0)\).
        Then \(x,y_0\in A_0\cap A\), with \(x\in C(A_0)\) and
        \(y_0\in C(A)\), so HARP yields \(x\in C(A)\).
      \end{itemize}
    \end{itemize}
  \end{itemize}
}



\rmkb{
\begin{enumerate}

\item
  HARP does the entire work in the ``if'' direction: it is precisely
  the axiom that endows the revealed preference relation \(\pref_C\)
  with the transitivity needed to call it rational.
\item
  The rationalizability result requires observing the \emph{whole}
  choice function \(C\) — that is, (i) for any given choice set, all of
  the agent's optimal choices are observed, not just some of them, and
  (ii) the agent's optimal choices are observed across \emph{every}
  choice set. Real data is rarely so generous:

  \begin{itemize}
  \item
    In consumer-choice problems, the menus \(A\) we see are typically
    budget sets indexed by prices and income — a strict subcollection
    of \(\B\), not all of \(\B\).
  \item
    The first form of incompleteness (only seeing some of the optimal
    choices from a given menu) is essentially intractable. For the
    second (only seeing some menus), other revealed-preference axioms
    have been developed. The \emph{Weak Axiom of Revealed Preference
    (WARP)} is the natural counterpart of HARP for choice rules
    restricted to budget sets and required to be single-valued (so the
    HARP conclusion sharpens to \(x = y\)). WARP is necessary for
    rationalizability but not sufficient on budget data alone; the
    \emph{Generalized Axiom of Revealed Preference (GARP)} closes the
    gap and is both necessary and sufficient.
  \end{itemize}
\end{enumerate}
}

